\begin{table}[tb]
\centering\footnotesize\setlength{\tabcolsep}{3pt}
\caption{保存入力の全６条件．各３予定生成．公開一致は既知使用量の比較可能対だけの一致数である．}
\begin{tabular}{lrrrr}
入力 & 試行 & 既知 & 完全非空 & 公開一致 \\ \hline
A & 3/3 & 3/3 & 3/3 & 3/3 \\
B & 3/3 & 3/3 & 3/3 & 3/3 \\
C & 3/3 & 3/3 & 3/3 & 3/3 \\
D & 3/3 & 3/3 & 3/3 & 3/3 \\
N & 3/3 & 3/3 & 3/3 & 3/3 \\
X & 3/3 & 3/3 & 3/3 & 3/3 \\
\hline\end{tabular}
\end{table}
公開出力はA/Bで0/9対，C/DとA/Nで各9/9対が一致した．全18生成は既知422,475論理token，使用量不明０で，推論器稼働は4,386.052秒だった．全15入力間比較と返却messageの不透明hashを保持し，内部推論の意味内容は読まない．
終了後の構造点検では，A/Bの差はJSON整形とコード位置情報を除いてもASTに残った．Aにはログインと曲・アルバム・プレイリスト一覧取得，Bにはログインと曲一覧取得の構文があった．これらは実行回数ではない．A/B/N/Xの12件は引数形式とPython構文の検査を通ったが，C/Dの６件は本文だけで構造化されたツール提出がなく，有効な操作の一致とは数えない．値・名前・文順は保持し，追加生成・コード実行・API・評価なしで全18件を点検した．出力の一致を実行の修復や課題品質へ読み替えない．
